\documentclass[final, letterpaper, 12pt]{article}
%%%%%%%%%%%%%%%%% Idioma y fuentes %%%%%%%%%
\usepackage[spanish,es-tabla]{babel}
\usepackage{fontspec} % para usar fuentes del sistema
\setmainfont{Times New Roman}[
  BoldFont={Times New Roman Bold}, 
  ItalicFont={Times New Roman Italic}, 
  BoldItalicFont={Times New Roman Bold Italic}]
\usepackage{amsmath, amsbsy}
\usepackage{unicode-math}
\setmathfont{Cambria Math}
\usepackage{microtype}
%%%%%%%%%%%%%%%%% Paquetes útiles %%%%%%%%%%
\usepackage{enumitem}
\usepackage[hypcap]{caption}
\captionsetup{font={small,it}, labelfont={small,it},}
\usepackage{graphicx}
\usepackage{multirow}
\usepackage{mathrsfs}
\let\eth\relax %arreglar error linea 20 de esta pagina

\usepackage{listings}
\usepackage{color}
\usepackage{longtable}
\usepackage{fancyhdr}
\usepackage{import}
\usepackage{float}
\usepackage{wrapfig}
\usepackage[margin=0.5in, foot=30pt, head=10pt, bottom=1.3in]{geometry}
\usepackage{hyphenat}
\usepackage[document]{ragged2e}
\usepackage{array}
\usepackage{epstopdf}
\usepackage{imakeidx}
\usepackage[bookmarks=true, hidelinks]{hyperref}
\usepackage{bookmark}
\usepackage{titlesec}
\usepackage{siunitx}
%\usepackage{gensymb}
\usepackage{setspace}
\usepackage{comment}
\usepackage{amsfonts}
\usepackage{makecell}
\usepackage{float}


%\usepackage{subfigure}
\usepackage{subcaption}
\usepackage{placeins}
%%%%%%%%%%%%%%%%% Configuración listings %%%%%%%%%%%%%%%
\definecolor{dkgreen}{rgb}{0,0.6,0}
\definecolor{darkblue}{rgb}{0,0,0.6}
\definecolor{darkgray}{rgb}{0.3,0.3,0.3}
\definecolor{lightgray}{rgb}{0.98,0.98,0.98}
\definecolor{mauve}{rgb}{0.58,0,0.82}
\lstset{frame=tb,
   language=Matlab,
   aboveskip=3mm,
   belowskip=3mm,
   showstringspaces=false,
   columns=flexible,
   basicstyle={\small\ttfamily},
   numbers=none,
   numberstyle=\tiny\color{gray},
   keywordstyle=\color{blue},
   commentstyle=\color{dkgreen},
   stringstyle=\color{mauve},
   breaklines=true,
   breakatwhitespace=true,
   tabsize=3}
%%%%%%%%%%%%%%%%% Encabezado y pie %%%%%%%%%%%%%%%
\geometry{left=2cm, right=2cm, top=2.5cm, bottom=3cm}
\pagestyle{fancy}
\fancyhead{}
\fancyfoot{}
\fancyfoot[L]{Grupo N°3}
\fancyfoot[R]{\thepage}
\fancyhead[R]{ELI327 - Laboratorio de máquinas eléctricas }
\fancyhead[L]{DIE, UTFSM}
\renewcommand{\footrulewidth}{0.9pt}
\setlength{\headheight}{0.09\linewidth}
%%%%%%%%%%%%%%%%% Documento %%%%%%%%%%%%%%%%%%%%%%
\begin{document}
\justifying

\begin{titlepage}
    \begin{center}
        \vspace*{0.8cm}
        
        \includegraphics[width=0.8\textwidth]{Logo-DIE-vertical-con-fondo.jpg}\\
        \vspace{0.9cm}
        \textbf{{\fontsize{14}{16}\selectfont \LARGE{Laboratorio de máquinas eléctricas}}}\\
        \vspace{0.65cm}
        \headrule
        \vspace{0.4cm}
        \textbf{{\fontsize{14}{16}\selectfont \LARGE{Protocolo 1 - Motor asincrónico sincronizado y estudio de asimetrías en rotor}}}\\
        \textbf{\uppercase{Grupo N$^\circ$3\\}}
        \vspace{0.4cm}
        \headrule
        \vspace{1.4cm}
Autores:\\
        Joaquín Vargas, ROL: 202023004-8\\
        \vspace{0.7cm}
        Profesor: Maximiliano Trujillo\\
        \vspace{1cm}
        Valparaíso, Chile\\
        \today
        
\end{center}       
\end{titlepage}



\newpage
\tableofcontents 
\newpage
\setlength{\parindent}{0pt} % sin sangría

\section{Objetivos}
\input{Desarrollo/Objetivos}
\newpage
\section{Teoría}
\input{Desarrollo/Teoria}
\newpage
\section{Método y desarrollo}
\input{Desarrollo/Metodo y desarrollo}
\newpage
\section{Ensayo}
\input{Desarrollo/Ensayo}
\newpage
\section{Crítica y comentarios}
\input{Desarrollo/Critica y comentarios}
\newpage
\section{Bibliografía}
\addcontentsline{toc}{section}{Referencias}
\begin{thebibliography}{100}
    \bibitem{primero} Este es un ejemplo
\end{thebibliography}
\newpage
\section{Apéndices}
\input{Desarrollo/Apendice}

\end{document}